\documentclass[a4paper,12pt]{ctexart} % 使用 ctexart 文档类，原生支持中文

% ==================== 宏包引入 ====================
\usepackage{geometry}       % 页面边距设置
\usepackage{graphicx}       % 插入图片
\usepackage{float}          % 图片位置强制固定
\usepackage{booktabs}       % 三线表
\usepackage{amsmath}        % 数学公式
\usepackage{listings}       % 代码块
\usepackage{xcolor}         % 颜色支持
\usepackage{hyperref}       % 目录跳转和超链接
\usepackage{caption}        % 图表标题格式
\usepackage{subcaption}

% ==================== 页面设置 ====================
\geometry{left=2.5cm, right=2.5cm, top=2.5cm, bottom=2.5cm}

% ==================== 代码高亮设置 ====================
\definecolor{codegreen}{rgb}{0,0.6,0}
\definecolor{codegray}{rgb}{0.5,0.5,0.5}
\definecolor{codepurple}{rgb}{0.58,0,0.82}
\definecolor{backcolour}{rgb}{0.95,0.95,0.92}

\lstdefinestyle{mystyle}{
    backgroundcolor=\color{backcolour},   
    commentstyle=\color{codegreen},
    keywordstyle=\color{magenta},
    numberstyle=\tiny\color{codegray},
    stringstyle=\color{codepurple},
    basicstyle=\ttfamily\footnotesize,
    breakatwhitespace=false,         
    breaklines=true,
    postbreak=\mbox{\textcolor{red}{$\hookrightarrow$}\space},
    captionpos=b,                    
    keepspaces=true,                 
    numbers=left,                    
    numbersep=5pt,                  
    showspaces=false,                
    showstringspaces=false,
    showtabs=false,                  
    tabsize=2,
    columns=flexible
}
\lstset{style=mystyle, language=Python}

% ==================== 文档信息 ====================
\title{\textbf{大模型课程第一次作业}}
\author{姓名：张宗桪 \\ 学号：2023K8009991013 \\ 专业：人工智能}
\date{\today}

% ==================== 正文开始 ====================
\begin{document}

\pagestyle{plain} % 页面顶部不显示标注，仅在底部显示页码

% 封面
\maketitle
\thispagestyle{empty} % 封面不显示页码
\newpage

% 目录
\tableofcontents
\newpage
\section{实验设计}


\subsection{数据集概况}
本次实验设置了两款数据集如下

\begin{table}[H]
    \centering
    \caption{实验所用数据集说明}
    \begin{tabular}{p{0.2\textwidth} p{0.72\textwidth}}
        \toprule
        数据集名称 & 数据内容 \\
        \midrule
        纯英文数据集 & 课程提供的英文新闻评论数据。 \\
        中英混合数据集 & 一半为新华日报官方标注数据集，另一半为上述英文新闻评论数据。 \\
        \bottomrule
    \end{tabular}
\end{table}

\subsection{核心超参数设置}

\begin{table}[H]
    \centering
    \caption{BPE 训练核心超参数设置}
    \begin{tabular}{p{0.3\textwidth} p{0.62\textwidth}}
        \toprule
        参数名称 & 参数设置 \\
        \midrule
        分词模型 & BPE \\
        词表大） & 6400 与 16000 \\
        最小词频 & 2 \\
        特殊符号 & \texttt{<|endoftext|>}、\texttt{<|im\_start|>}、\texttt{<|im\_end|>} \\
        未知词标记 & \texttt{<|endoftext|>} \\

        \bottomrule
    \end{tabular}
\end{table}


\section{实验结果}


\subsection{训练结果}


\subsection{测试结果}


\section{结果分析}
% -------------------- 第一部分 --------------------

\end{document}